\documentclass[11pt,a4paper]{ctexart}
\usepackage[teacher]{../../../common/ipara}
\tcbuselibrary{breakable}
\usepackage{multicol}

\usepackage{pgfplots}
\pgfplotsset{compat=1.18}
\usetikzlibrary{arrows.meta,positioning,decorations.pathreplacing}

\setlist[itemize]{leftmargin=*}
\setlist[enumerate]{leftmargin=*}

\pagestyle{fancy}
\fancyhf{}
\lhead{\small 高一/高三数学专题：函数同构与指对同构专题讲义}
\rhead{\small 教师版（普通同构与指对变形技术）}
\cfoot{\small 第 \thepage\ 页\quad 共 \pageref{LastPage} 页}
\renewcommand{\headrulewidth}{0.4pt}
\renewcommand{\footrulewidth}{0pt}


\renewcommand{\teacherproblem}[2]{%
  \teachquestionheader{例题 #1\quad #2}
}

\begin{document}

\begin{center}
  {\zihao{2}\bfseries 函数观点下的代数变形：函数同构与指对同构}\\[3mm]
  {\Large \bfseries 教师备课与授课专题讲义}\\[2mm]
  {\large 适用：高一函数综合提升 / 高三一轮与导数大题专题 / 拔高课次}\\[2mm]
  \rule{0.85\linewidth}{0.5pt}
\end{center}

\begin{teachbox}{专题设计定位}
本专题重在引导学生突破高考中档与拔高题中最常见的\textbf{“同构变形与函数构造”}技巧。
同构不仅是解决代数式繁琐计算的“神技”，更是\textbf{函数观点（代数式形式与对应函数结构的一致性）}的核心体现。
本讲分成两层：
\begin{enumerate}
  \item \textbf{第一层：普通代数同构}。通过移项使等式或不等式两边结构完全一致，进而构造辅助函数，借助单调性解决变量关系。
  \item \textbf{第二层：指对同构（本课重点）}。利用指数与对数的恒等变换（如 $x=e^{\ln x}$, $y=\ln(e^y)$），将表面上完全不同构的“幂指混合”与“幂对混合”式子，强行化为同一函数结构，结合导数与单调性分析。
\end{enumerate}
\end{teachbox}

\newpage

\section{第一部分：普通代数同构与函数构造}

\subsection{1. 普通同构的判断与构造流程}
判断一个式子是否同构，其核心是寻找\textbf{形式与变量的归位}。
如果我们能通过移项或代数化简，将关于变量 $x_1$ 的部分全部集中到等号左侧，将关于变量 $x_2$ 的部分集中到等号右侧，且左右两边的代数表达结构完全一致：
\[
F(x_1) = F(x_2) \quad \text{或} \quad F(x_1) \ge F(x_2)
\]
那么，这个共同的结构 $F(t)$ 就是我们要构造的辅助函数。
解题流程总结为：
\begin{enumerate}
  \item \textbf{归项归位}：将同一变量相关的项放在同侧。
  \item \textbf{观察结构}：把变量用统一符号“$\square$”代替，检查两边结构是否完全对齐。
  \item \textbf{构造函数}：将“$\square$”抽象为自变量 $t$，写出辅助函数 $g(t)$ 及其定义域。
  \item \textbf{研究单调}：求导或利用基本函数性质判定 $g(t)$ 在特定区间上的单调性，利用单调性脱去外层函数 $g$，转化为变量的大小关系。
\end{enumerate}

\begin{exbox}{例 1.1}
已知函数 $f(x)$ 满足对任意两个不同的实数 $x_1, x_2$ 均有：
\[
f(x_1) - 2x_1 \ge f(x_2) - 2x_2 \quad (x_1 > x_2)
\]
试说明其几何意义，并写出对应构造函数的性质。
\end{exbox}

\begin{paracol}{2}
\teachblock{标准解答与讲解主线}{
  \textbf{解题思路与构造：}\\
  由于 $x_1 > x_2$，我们观察式子，将同类变量集中在同侧。原不等式已经分居两侧：
  \[
  f(x_1) - 2x_1 \ge f(x_2) - 2x_2
  \]
  观察左边和右边的结构，它们都是：
  \[
  f(\text{变量}) - 2\text{变量}
  \]
  因此，我们可以构造辅助函数：
  \[
  g(x) = f(x) - 2x
  \]
  那么原式即可翻译为：
  当 $x_1 > x_2$ 时，恒有：
  \[
  g(x_1) \ge g(x_2)
  \]
  根据单调性的定义，这说明构造出的辅助函数 $g(x) = f(x) - 2x$ 在其定义域上是\textbf{单调递增}的。
}
\switchcolumn
\teachblock{易错分析与教学提示}{
  \textbf{学生常见误区：}\\
  学生容易直接去对 $f(x)$ 猜想函数解析式，或者不知道如何对 $f(x)$ 与 $2x$ 进行结合。
  
  \textbf{启发与提问口径：}\\
  \begin{itemize}
    \item “我们看左边是关于 $x_1$ 的式子，右边是关于 $x_2$ 的式子。如果我把 $x_1$ 和 $x_2$ 用括号括起来，两边的代数长相是不是完全一样？这种长相就可以抽象成一个统一的函数 $g(x) = f(x) - 2x$。”
  \end{itemize}
}
\end{paracol}

\begin{exbox}{例 1.2}
若对于任意实数 $x_1, x_2 \in [e, +\infty)$，恒有 $\dfrac{\ln x_1}{x_1} \le \dfrac{\ln x_2}{x_2}$，试比较 $x_1$ 与 $x_2$ 的大小。
\end{exbox}

\begin{paracol}{2}
\teachblock{标准解答与讲解主线}{
  \textbf{解题步骤：}\\
  观察两边结构，均呈 $\dfrac{\ln(\text{变量})}{\text{变量}}$ 的形式。
  
  因此，我们构造辅助函数：
  \[
  g(t) = \frac{\ln t}{t} \quad (t \ge e)
  \]
  那么原不等式即可写为：
  \[
  g(x_1) \le g(x_2)
  \]
  接下来我们研究 $g(t)$ 的单调性。对 $g(t)$ 求导：
  \[
  g'(t) = \frac{(\ln t)' \cdot t - \ln t \cdot t'}{t^2} = \frac{1 - \ln t}{t^2}
  \]
  因为自变量 $t \ge e$，所以 $\ln t \ge 1$，从而可得：
  \[
  1 - \ln t \le 0 \implies g'(t) \le 0
  \]
  故函数 $g(t)$ 在 $[e, +\infty)$ 上是\textbf{严格单调递减}的。
  
  根据单调递减函数的性质，有：
  \[
  g(x_1) \le g(x_2) \implies \boxed{x_1 \ge x_2}
  \]
}
\switchcolumn
\teachblock{易错分析与教学提示}{
  \textbf{学生常见误区：}\\
  学生容易忽略对 $g(t)$ 单调递减时的方向变号。认为 $g(x_1) \le g(x_2)$ 必然推出 $x_1 \le x_2$。
  
  \textbf{启发与提问口径：}\\
  \begin{itemize}
    \item “我们算出了 $g(t)$ 在 $[e, +\infty)$ 上是递减的。递减函数就像坐滑梯，自变量越大，函数值越小。现在 $g(x_1)$ 的值比 $g(x_2)$ 还要低，说明 $x_1$ 和 $x_2$ 谁在更右边（更大）？显然是 $x_1 \ge x_2$。”
  \end{itemize}
}
\end{paracol}

\newpage

\section{第二部分：指对同构（本课重点）}

\subsection{1. 指对同构的核心思想与变形基础}
指对同构是处理“指数与对数混合式”的强效工具。这类题目表面上左右两边形式截然不同（一边是指数，一边是对数），但通过指数与对数的恒等变形公式，它们可以化归为同一种代数结构。

\begin{keybox}{指对同构 4 大基本恒等式}
任何正实数或实数都可以通过指对互化重构其形式：
\begin{itemize}
  \item $y = \ln(e^y)$ \quad（将代数式改写为对数形式）
  \item $y = \log_a(a^y)$
  \item $x = e^{\ln x} \quad (x > 0)$ \quad（将代数式改写为指数底数形式）
  \item $x = a^{\log_a x} \quad (x > 0)$
\end{itemize}
\end{keybox}

\subsection{2. 指对同构的经典对齐结构}
在实际题目中，我们通常将“幂指”结构与“幂对”结构进行对称转换：
\begin{center}
\renewcommand{\arraystretch}{1.3}
\begin{tabularx}{0.95\linewidth}{c c X}
\toprule
结构类型 & 原始形式组 & 同构对齐后的统一结构 \\
\midrule
乘积型 & $x \ln x \quad \text{与} \quad y e^y$ & 变形为：$e^{\ln x}\ln x = y e^y$ （统一为 $t e^t$） \\
 & & 或：$x \ln x = e^y \ln(e^y)$ （统一为 $t \ln t$） \\
\midrule
和差型 & $x + 3^x \quad \text{与} \quad y + \log_3 y$ & 变形为：$\log_3(3^x) + 3^x = \log_3 y + y$ （统一为 $\log_3 t + t$） \\
\midrule
差值型 & $x - \ln x \quad \text{与} \quad e^y - y$ & 变形为：$e^{\ln x} - \ln x = e^y - y$ （统一为 $e^t - t$） \\
\midrule
商式型 & $x / \ln x \quad \text{与} \quad e^y / y$ & 变形为：$e^{\ln x} / \ln x = e^y / y$ （统一为 $e^t / t$） \\
\bottomrule
\end{tabularx}
\end{center}

\newpage

\begin{exbox}{例 2.1（由两方程根建立关系）}
已知实数 $x_1$ 是方程 $\log_2 x + x = 7$ 的根，实数 $x_2$ 是方程 $2^x + x = 7$ 的根，求 $x_1 + x_2$ 的值。
\end{exbox}

\begin{paracol}{2}
\teachblock{标准解答与讲解主线}{
  \textbf{步骤一：代入根建立等式关系}\\
  由题意，将根代入方程：
  \[
  \log_2 x_1 + x_1 = 7, \qquad 2^{x_2} + x_2 = 7
  \]
  从而得到等量关系：
  \[
  \log_2 x_1 + x_1 = 2^{x_2} + x_2
  \]
  
  \textbf{步骤二：利用指对恒等式进行同构对齐}\\
  左侧是“对数 + 幂”，右侧是“指数 + 幂”，结构不对齐。
  我们将右侧的自变量 $x_2$ 改写为对数形式：$x_2 = \log_2(2^{x_2})$。
  代入等式右侧：
  \[
  \log_2 x_1 + x_1 = \log_2(2^{x_2}) + 2^{x_2}
  \]
  
  \textbf{步骤三：构造函数并判定单调性}\\
  两边结构完全统一为 $\log_2 t + t$。
  构造辅助函数：
  \[
  f(t) = \log_2 t + t \quad (t > 0)
  \]
  由于 $y=\log_2 t$ 和 $y=t$ 在 $(0, +\infty)$ 上均是严格单调递增的，故它们的和函数 $f(t)$ 也是严格单调递增的。
  
  \textbf{步骤四：脱去函数符号，求变量关系}\\
  因为 $f(x_1) = f(2^{x_2})$ 且 $f(t)$ 单调，所以：
  \[
  x_1 = 2^{x_2}
  \]
  
  \textbf{步骤五：结合原方程求解}\\
  代回方程 $2^{x_2} + x_2 = 7$，即可得：
  \[
  x_1 + x_2 = \boxed{7}
  \]
}
\switchcolumn
\teachblock{易错分析与教学提示}{
  \textbf{学生常见误区：}\\
  学生容易试图强行求解这两个超越方程。由于涉及指数与对数，根本无法用初等方法求出 $x_1$ 和 $x_2$ 的具体数值。
  
  \textbf{启发与提问口径：}\\
  \begin{itemize}
    \item “我们能单独把 $x_1$ 或者 $x_2$ 算出来吗？算不出来，因为它们是超越方程。但这道题问的是它们的和 $x_1+x_2$。这启发我们：不要试图求单体，而是通过‘同构’建立 $x_1$ 和 $x_2$ 之间的关联结构！”
  \end{itemize}
}
\end{paracol}

\newpage

\begin{exbox}{例 2.2（结构需先调整的同构）}
已知实数 $x, y$ 满足方程组：
\[
\begin{cases}
  x + 3^{3x-1} = 2 \\
  9y + \log_3 y = 4
\end{cases}
\]
求 $x + 3y$ 的值。
\end{exbox}

\begin{paracol}{2}
\teachblock{标准解答与讲解主线}{
  \textbf{步骤一：两端调整，对齐内部整体}\\
  第一式中指数为 $3x-1$，旁边为 $x$。
  为配凑出 $3x$，我们将第一式两边同乘以 3：
  \[
  3x + 3 \cdot 3^{3x-1} = 6 \implies 3x + 3^{3x} = 6
  \]
  第二式中对数为 $\log_3 y$，旁边为 $9y$。
  利用对数公式变形：$\log_3(9y) = \log_3 y + 2 \implies \log_3 y = \log_3(9y) - 2$。
  代入第二式：
  \[
  9y + \log_3(9y) - 2 = 4 \implies 9y + \log_3(9y) = 6
  \]
  现在，两方程的常数端均化为 6，我们得到联立等式：
  \[
  3x + 3^{3x} = 9y + \log_3(9y)
  \]
  
  \textbf{步骤二：引入恒等变形公式}\\
  我们将左侧的 $3x$ 改写为对数形式：$3x = \log_3(3^{3x})$。
  代入式子中可得：
  \[
  \log_3(3^{3x}) + 3^{3x} = \log_3(9y) + 9y
  \]
  
  \textbf{步骤三：构造函数与脱去符号}\\
  构造辅助函数 $f(t) = \log_3 t + t$ ($t > 0$)。
  易知 $f(t)$ 在 $(0, +\infty)$ 上严格单调递增。
  因为 $f(3^{3x}) = f(9y)$，故：
  \[
  3^{3x} = 9y
  \]
  
  \textbf{步骤四：代回原式求值}\\
  代回变形后的第一式 $3x + 3^{3x} = 6$，有：
  \[
  3x + 9y = 6 \implies x + 3y = \boxed{2}
  \]
}
\switchcolumn
\teachblock{易错分析与教学提示}{
  \textbf{学生常见误区：}\\
  学生看到第一式中的 $x$ 和第二式中的 $9y$ 结构不匹配，往往不知道该如何下手，卡在起步阶段。
  
  \textbf{启发与提问口径：}\\
  \begin{itemize}
    \item “我们看第一式里的指数上有 $3x$，所以旁边孤零零的 $x$ 应该主动乘个 3。同理，第二式对数真数里是 $y$，旁边是 $9y$。我们把对数强行凑出 $9y$，就用 $\log_3(9y)$。这种‘给什么凑什么’的整体法是同构的关键前哨战。”
  \end{itemize}
}
\end{paracol}

\newpage

\begin{exbox}{例 2.3（同构后需严格检查单调区间的题目）}
若正实数 $x_1, x_2$ 分别满足方程：
\[
x_1 \ln x_1 = 1, \qquad x_2 e^{x_2} = 1
\]
求 $x_1 \cdot x_2$ 的值。
\end{exbox}

\begin{paracol}{2}
\teachblock{标准解答与讲解主线}{
  \textbf{步骤一：指对变换统一结构}\\
  因为 $x_1 > 0$，我们将第一个方程中的自变量 $x_1$ 变换为指数形式：$x_1 = e^{\ln x_1}$。
  代入第一式：
  \[
  \ln x_1 \cdot e^{\ln x_1} = 1
  \]
  第二式原形为：
  \[
  x_2 e^{x_2} = 1
  \]
  
  \textbf{步骤二：构造函数并求导分析}\\
  我们得到两端结构一致的等式：$\ln x_1 \cdot e^{\ln x_1} = x_2 e^{x_2}$。
  构造辅助函数：
  \[
  f(t) = t e^t
  \]
  从而有 $f(\ln x_1) = f(x_2)$。对 $f(t)$ 求导：
  \[
  f'(t) = e^t + t e^t = (t+1)e^t
  \]
  令 $f'(t) > 0 \implies t > -1$；令 $f'(t) < 0 \implies t < -1$。
  故函数 $f(t)$ 在 $(-\infty, -1)$ 上单调递减，在 $(-1, +\infty)$ 上单调递增。
  
  \textbf{步骤三：确定自变量范围，锁定单调区间}\\
  因为函数 $f(t)$ 不是全区间单调的，我们必须严格证明两个自变量 $\ln x_1$ 与 $x_2$ 落在同一个单调区间内：
  \begin{itemize}
    \item 对于 $x_2$：因为 $x_2 e^{x_2} = 1 > 0$ 且 $e^{x_2} > 0$，所以必然有 $x_2 > 0$。
    \item 对于 $\ln x_1$：因为 $x_1 \ln x_1 = 1 > 0$ 且 $x_1 > 0$，所以必然有 $\ln x_1 > 0$。
  \end{itemize}
  因为两者均大于 0，显然均落在区间 $(0, +\infty) \subset (-1, +\infty)$ 内。
  在此递增区间上，由 $f(\ln x_1) = f(x_2)$ 才能安全推出：
  \[
  \ln x_1 = x_2 \implies x_1 = e^{x_2}
  \]
  
  \textbf{步骤四：求值}\\
  \[
  x_1 \cdot x_2 = e^{x_2} \cdot x_2 = 1
  \]
  即：$\boxed{x_1 \cdot x_2 = 1}$。
}
\switchcolumn
\teachblock{易错分析与教学提示}{
  \textbf{学生常见误区：}\\
  学生在由 $f(\ln x_1) = f(x_2)$ 推出 $\ln x_1 = x_2$ 时，几乎全员都会漏掉对“单调区间”的限制论证。如果辅助函数在实数集上不单调（如二次函数），$f(a)=f(b)$ 并不代表 $a=b$！
  
  \textbf{启发与提问口径：}\\
  \begin{itemize}
    \item “构造出 $f(t)=te^t$ 之后，我们算它的对称折返点是 $t=-1$。在左边递减，右边递增。我们必须向判卷老师证明：$\ln x_1$ 和 $x_2$ 都大于 0，也就是说它们都在同一个只增不减的滑梯上。这样它们对应的自变量才能完全相等。”
  \end{itemize}
}
\end{paracol}

\newpage

\section{第三部分：解题决策与思维流程总结}

\subsection{1. 备课研讨：如何看穿“被打乱”的同构式？——增长率排序定位法}

在实际命题中，出题人为了增加难度，往往会把等式或不等式的各项顺序完全打乱（例如将所有项移到一侧），使学生很难一眼看出左右两边到底是不是同构，甚至不知道该如何分配项。

此时，教师可以引导学生使用\textbf{“增长率分层排序定位法”}。

根据基本初等函数在自变量趋于正无穷时的增长速率（即导数的数量级），我们有如下的分层排位：
\[
\text{指数级 } (e^t, a^t) \gg \text{幂级 / 线性级 } (t, mt+n) \gg \text{对数级 } (\ln t, \log_a t)
\]
当遇到打乱的指对混合式时，可以通过以下步骤理清结构：
\begin{enumerate}
  \item \textbf{增长率归类}：将等式中的各项划分为指数级、线性级、对数级三个档次。
  \item \textbf{对齐归位}：将高增长率与低增长率的项在两侧进行配对。例如，左边配对“指数 $\times$ 线性”，则右边对应的目标应当是“线性 $\times$ 对数”（因为对数进行指对互化后可以提升一个增长等级）。
  \item \textbf{转换验证}：使用 $y = e^{\ln y}$ 将较低等级的项转化为较高等级的项，或者使用 $y = \ln(e^y)$ 降级，看两边是否能够实现增长率级别的完全对称。
\end{enumerate}

\begin{teachbox}{备课微实例}
例如：方程 $x e^x = y \ln y$ 中，$x, y$ 均为正数，试通过增长率排序法找出其同构结构。
\begin{itemize}
  \item \textbf{分析}：左边是 $x e^x$（线性级 $\times$ 指数级）；右边是 $y \ln y$（线性级 $\times$ 对数级）。
  \item \textbf{对齐}：左边结构为 $t e^t$（自变量与其对应指数相乘）。为了在右边也凑出类似的“指数 $\times$ 自变量”结构，我们将较低增长等级的线性级 $y$ 升级为指数级 $e^{\ln y}$。
  \item \textbf{变形}：
    \[
    y \ln y = e^{\ln y} \ln y = \ln y \cdot e^{\ln y}
    \]
    此时，原式化为：
    \[
    x e^x = \ln y \cdot e^{\ln y}
    \]
    左右两边完全对齐为 $t e^t$ 的结构！这正是通过增长速率分层对比，指导我们精准将 $y$ 升格为 $e^{\ln y}$ 的决策逻辑。
\end{itemize}
\end{teachbox}

\subsection{2. 函数同构问题“大脑扫描决策 5 步流程”}

\begin{center}
\begin{tcolorbox}[colback=white, colframe=black, boxrule=0.5pt, arc=0pt, left=8pt, right=8pt, top=8pt, bottom=8pt]
\textbf{【同构与指对同构解题决策模板】}
\begin{enumerate}[label=第\arabic*步：, leftmargin=3.5em]
  \item \textbf{整体对齐检验}：观察方程或不等式中是否出现“指数与幂”或“对数与幂”的配对。若有，则进入指对同构准备。
  \item \textbf{对齐系数整体}：优先把相同底数下的整体调整一致（如看到 $3^{3x}$，就把旁边的自变量乘系数变成 $3x$；看到 $\log_3(9y)$，就让旁边出现 $9y$）。
  \item \textbf{套用指对恒等式}：根据凑配方向，使用公式 $x = e^{\ln x}$ 或 $y = \ln(e^y)$ 强行对齐两边长相。
  \item \textbf{构造函数与求导}：抽象出形式一致的辅助函数 $f(t)$。若定义域内不完全单调，必须单独求导判定其极值点。
  \item \textbf{单调区间锁定与脱壳}：结合原方程正负条件，说明两个变量同属一个单调区间。脱去函数外壳 $f$，求出变量等量关系并代回计算。
\end{enumerate}
\end{tcolorbox}
\end{center}

\end{document}
